% ECONOMICS_PROBLEM: {"id":"C73-17","title":"Maxmin equality with the concavified nonrevealing value","jel_code":"C73","source_id":"OP-0168"}
% FIRST_POSTED: 2026-10-09
% ATTEMPT_STATUS: unresolved
% ENTRY_KIND: research_attempt
\documentclass[11pt]{article}
\usepackage{amsmath,amssymb,amsthm,geometry,hyperref}
\geometry{margin=1in}
\newtheorem{theorem}{Theorem}
\newtheorem{proposition}[theorem]{Proposition}
\newtheorem{lemma}[theorem]{Lemma}
\newtheorem{corollary}[theorem]{Corollary}
\theoremstyle{definition}
\newtheorem{example}[theorem]{Example}
\newtheorem{remark}[theorem]{Remark}
\newcommand{\R}{\mathbb R}
\newcommand{\N}{\mathbb N}
\DeclareMathOperator{\NE}{NE}
\DeclareMathOperator{\cav}{cav}
\DeclareMathOperator{\supp}{supp}
\title{Maxmin equality with the concavified nonrevealing value\\ GPT 6 Astra Ultra}
\author{Research attempt}
\date{9 October 2026}
\begin{document}
\maketitle

\section*{Target and contribution}
The target asks for structural conditions on a bounded Borel payoff functional $f(k,\omega)$ ensuring $\underline v_f(p)=(\cav u_f)(p)$ for every prior. Here $\underline v_f$ is the informed maxmin and $u_f$ the nonrevealing value. The known tail-measurable theorem is already outside the open portion of the question \cite{tail}. This attempt proves the general direction of the inequality carefully, then derives a complete criterion for a small first-stage family and a sufficient condition for general finite first-stage games. These are elementary restricted results, not a new general characterization.

\section*{Why a counterexample can only lie above concavification}
\begin{lemma}[Concavity of the informed maxmin]
For every finite state set, finite action alphabets, and bounded Borel payoff $f$, the informed maxmin $\underline v_f$ is concave on the prior simplex. Moreover
\[
 \underline v_f(p)\ge(\cav u_f)(p)
 \quad\text{for every prior }p.
\]
\end{lemma}
\begin{proof}
Take a finite splitting $p=\sum_r\lambda_r q_r$. For each $r$, choose a behavioral informed strategy $\sigma^r$ guaranteeing at least $\underline v_f(q_r)-\epsilon$ at prior $q_r$. For $p(k)>0$, privately choose an index $r$ with probability $\lambda_rq_r(k)/p(k)$, conditional on state $k$, and then follow $\sigma^r$. Conditional on the chosen index the state law is $q_r$. An uninformed opponent still uses some strategy depending only on the public history. Thus the conditional expected payoff is at least $\underline v_f(q_r)-\epsilon$, giving the weighted guarantee.

To justify this construction within the stated behavioral-strategy convention, rather than implicitly granting an extra remembered signal, define at public history $h=(i_1,j_1,\ldots,i_t,j_t)$ the weights
\[
 w_r(k,h)=\frac{\lambda_rq_r(k)}{p(k)}
             \prod_{s=1}^t\sigma^r(i_s\mid k,h_{s-1}).
\]
At a history with positive total weight, play the normalized weighted average of the next-action distributions $\sigma^r(\cdot\mid k,h)$. At a zero-weight history choose arbitrarily. Multiplying these conditional probabilities along a finite history telescopes the normalizing weights and produces exactly the mixture of the $\sigma^r$ likelihoods. The likelihood contributions of the opponent are common to every $r$ and cancel in the posterior weights. Hence this behavioral strategy induces the same probabilities on every finite cylinder as the private-index construction, against every opponent. The laws on the infinite product coincide, so bounded Borel payoffs also have equal expectations. States with $p(k)=0$ are irrelevant and may use arbitrary prescriptions. Letting $\epsilon\downarrow0$ proves concavity.

The informed player may ignore its state. For any state-independent strategy pair the expected payoff is the payoff in the complete-information functional $f_p$. Using an $\epsilon$-optimal maximizing strategy for that game gives $\underline v_f(p)\ge u_f(p)$. Apply the concavity just proved to every finite splitting of $p$, replace each $\underline v_f(q_r)$ by $u_f(q_r)$, and take the supremum defining $\cav u_f$.
\end{proof}

The inequality is known in the source; the proof is included to make the quantifiers and the admissibility of hidden randomization explicit. No claimed minimax equality for the incomplete-information Borel game is needed.

\section*{A fully solved first-stage parameter family}
Fix states $K=\{0,1\}$ and two actions for each player. Let $p$ be the probability of state $0$. For parameters $a,b\ge0$, let the payoff depend only on the first joint action and be given by the matrices
\[
 A_0=\begin{pmatrix}a&0\\0&0\end{pmatrix},\qquad
 A_1=\begin{pmatrix}0&0\\0&b\end{pmatrix}.
\]
Subsequent actions have no payoff effect, so these are bounded Borel functionals, generally not tail measurable.

\begin{theorem}[Exact values and criterion]
If $a,b>0$, then for every $p\in[0,1]$,
\[
 u_f(p)=\frac{abp(1-p)}{ap+b(1-p)},\qquad
 (\cav u_f)(p)=u_f(p),\qquad
 \underline v_f(p)=\min\{ap,b(1-p)\}.
\]
For every interior prior the inequality $\underline v_f(p)>(\cav u_f)(p)$ is strict. Over the entire parameter family $a,b\ge0$, equality for every prior holds if and only if $ab=0$.
\end{theorem}
\begin{proof}
Put $A=ap$, $B=b(1-p)$. The nonrevealing payoff matrix is diagonal with entries $A,B$. If the row player chooses row one with probability $r$, its payoffs against the two columns are $Ar$ and $B(1-r)$. For $A+B>0$, their minimum is maximized at $r=B/(A+B)$, with value $AB/(A+B)$. The column player can attain the same upper bound by choosing column one with probability $B/(A+B)$, because then both pure rows have payoff $AB/(A+B)$. This proves the value without any unproved minimax interchange. At an endpoint prior, one diagonal entry vanishes and the value is zero, as the displayed formula states.

For $a,b>0$, differentiation gives
\[
 u_f''(p)=-\frac{2a^2b^2}{[b+(a-b)p]^3}<0.
\]
Thus $u_f$ is concave on the closed interval, by its continuous endpoint extension, and its concavification is itself.

The informed player can choose row one in state $0$ and row two in state $1$. Against a column-one probability $y$ its expected payoff is $Ay+B(1-y)$, whose minimum over $y\in[0,1]$ is $\min\{A,B\}$. Conversely the column player may choose column one, guaranteeing payoff at most $A$ against every informed row strategy, or column two, guaranteeing at most $B$. This proves the informed value. If $A,B>0$, then $AB/(A+B)$ is strictly below both $A$ and $B$. At every interior prior they are both positive, proving the strict inequality. If $a=0$ or $b=0$, a column gives zero in both states against all rows; all payoffs are nonnegative, so all three quantities are identically zero. These cases prove the claimed if-and-only-if condition.
\end{proof}

For $a=b=1$ and $p=1/2$, the exact comparison is $\underline v_f=1/2$ and $\cav u_f=1/4$. This two-state, two-action calculation is a small adversarial check on any proposed extension from tail payoffs to arbitrary finite-prefix dependence.

\section*{A sufficient condition allowing prefix dependence}
\begin{proposition}[A common minimizing strategy across states]
Let $f(k,\omega)=A_k(i_1,j_1)$ depend only on the first action pair, with finite matrices $A_k$. Suppose there exist numbers $v_k$, row distributions $x^k$, and one column distribution $y^*$ such that for every state $k$,
\[
 A_k(x^k,j)\ge v_k\quad\text{for every pure column }j,
 \qquad
 A_k(i,y^*)\le v_k\quad\text{for every pure row }i.
\]
Then for every prior $p$,
\[
 \underline v_f(p)=(\cav u_f)(p)=\sum_k p(k)v_k.
\]
\end{proposition}
\begin{proof}
The informed player uses $x^k$ in state $k$ and guarantees the displayed affine quantity. The opponent's common strategy $y^*$ bounds the payoff of every informed strategy above by the same quantity. Thus the informed value equals it. In the nonrevealing game, playing $y^*$ gives $u_f(q)\le\sum_kq(k)v_k$ for every $q$. At each vertex $e_k$, the two assumed inequalities give $u_f(e_k)=v_k$. Every splitting therefore has weighted nonrevealing value at most $\sum_kp(k)v_k$, while the splitting into the vertices attains exactly this amount. The formula for the concavification follows.
\end{proof}

This condition is an explicit finite matrix certificate. It includes non-tail payoffs, so tail measurability is not necessary. It is not asserted necessary: the general characterization must account for other mechanisms as well.

\section*{Remaining gap}
The general lower bound, the exact first-stage parameter classification, and the common-minimizer certificate do not provide necessary and sufficient structural conditions for arbitrary bounded Borel functionals. In particular, concatenating a finite information-dependent block with a tail game may change both the available signals and the payoff. No invariance under such concatenation is assumed or proved.

\begin{thebibliography}{9}
\bibitem{tail} G. Bar Castellon Koltun, E. Lehrer, and E. Solan, \emph{The Maxmin Value of Repeated Games with Incomplete Information on One Side and Tail-Measurable Payoffs}, arXiv:2512.01581 (2025), Definitions 2.1--2.4, Theorem 2.5, the first part of its proof, and Section 5. \url{https://arxiv.org/pdf/2512.01581}. The primary PDF, its tail-measurable equality, and its general lower-bound discussion were checked on 9 October 2026. The results in this attempt are proved directly; no existence of an incomplete-information game value is inferred from the tail theorem.
\end{thebibliography}

\end{document}
